This study presented HS-AeroTS, a hierarchical framework that connects UAV telemetry anomaly detection with fault-domain diagnosis, interpretable telemetry evidence, and architecture-aware software inspection. The framework transforms multivariate PX4 telemetry into temporal-statistical representations, separates anomaly screening from conditional fault-domain diagnosis, and propagates class-specific TreeSHAP evidence from statistical features to telemetry channels, uORB topics, and firmware-specific PX4 module candidates.

The experiments show that the two predictive stages exhibit different generalization characteristics. Under leave-log-out evaluation, Stage 1 achieved an AUPRC of $0.6296\pm0.0023$, while the conditional four-domain classifier achieved a Macro-F1 of $0.9041\pm0.0029$. Comparisons across chronological, purged, and flight-isolated protocols indicate that anomaly screening is more sensitive to cross-flight variation than conditional fault-domain diagnosis. The hierarchical formulation should therefore be interpreted primarily as a diagnostic organization that separates detection, fault interpretation, and evidence generation rather than as a mechanism that guarantees higher complete five-class accuracy. Attribution-guided masking further confirmed that the identified TreeSHAP features contain prediction-relevant information, and the commit-specific uORB mapping provided a traceable mechanism for converting telemetry evidence into software-module inspection candidates.

Controlled replay provided a complementary source-level evaluation of this evidence chain. Common-support residual processing reduced paired healthy-run alarms from 11/12 to 3/12 while retaining 10/12 fault detections, with Top-3 and Top-5 module coverage of 10/12. These results demonstrate that matched telemetry evidence can narrow the software inspection space under controlled conditions, but persistent EKF2 false alarms, weak or delayed INAV responses, and the dependence on matched healthy references prevent the current approach from being interpreted as reliable end-to-end software fault localization. Overall, HS-AeroTS provides a reproducible bridge from real-flight telemetry diagnosis to architecture-aware software inspection. Future work should focus on more robust cross-flight anomaly screening, runtime-aware software provenance, and fully independent source-level validation across broader PX4 revisions and fault scenarios.
